\documentclass[x11names,svgnames, usepdftitle=false]{beamer}

\usetheme{metropolis}
\setbeamertemplate{caption}{\raggedright\insertcaption\par}
\setbeamercolor{block title}{fg=LightSkyBlue4}
\newenvironment<>{importantblock}[1]{%
  \begin{actionenv}#2%
    \def\insertblocktitle{#1}%
    \par%
    \mode<presentation>{%
      \setbeamercolor{block title}{%
        use=normal text,
        fg=normal text.fg,
        bg=normal text.bg!80!fg
      }
      \setbeamercolor{block body}{
        use={block title, normal text},
        bg=block title.bg!50!normal text.bg
      }
    }%
    \usebeamertemplate{block begin}}
  {\par\usebeamertemplate{block end}\end{actionenv}}
\newenvironment<>{decisionblock}[1][Statement of the problem]{%
  \begin{actionenv}#2%
    \def\insertblocktitle{#1}%
    \par%
    \mode<presentation>{%
      \setbeamercolor{block title}{%
        use=normal text,
        fg=normal text.fg,
        bg=normal text.bg!80!fg
      }
      \setbeamercolor{block body}{
        use={block title, normal text},
        bg=block title.bg!50!normal text.bg
      }
    }%
    \usebeamertemplate{block begin}}
  {\par\usebeamertemplate{block end}\end{actionenv}}
\usepackage{fontspec}
\usepackage[french, frenchkw]{algorithm2e}
\SetKwFor{Pc}{Pour chaque}{faire}{fin}
\SetKwProg{Fn}{Fonction}{:}{}
\usepackage{textcomp}
\usepackage{minted}
\usepackage{pdfpages}
\usepackage{qrcode}
\usepackage[francais]{babel}
\usepackage{xspace}
\usepackage[normalem]{ulem}
\usepackage[backend=biber,style=authoryear,url=false]{biblatex}
\addbibresource{Bibliography.bib} % Bibliography file
\usepackage{amsmath}
\usepackage{mathtools}
\usepackage{pifont}% http://ctan.org/pkg/pifont
\newcommand{\cmark}{\ding{51}}%
\newcommand{\xmark}{\ding{55}}%
\usepackage{marvosym}
\usepackage{tikz}
\usetikzlibrary{calc, positioning, arrows, automata, fadings, decorations.pathmorphing, shapes, patterns, tikzmark, fit, chains, snakes, shapes.multipart}

% ------------------------------------------------------------------
% Illustration de l'exécution des fonctions (plan / appel)
% Code couleur : violet = paramètres (les « trous » du plan)
%                orange = valeurs (arguments, valeur renvoyée)
% ------------------------------------------------------------------
\definecolor{paramcol}{RGB}{128, 0, 160}
\definecolor{valcol}{RGB}{214, 106, 0}
\definecolor{kwcol}{RGB}{0, 110, 0}
\definecolor{blueprintfill}{RGB}{225, 238, 250}
\definecolor{blueprintline}{RGB}{70, 130, 180}

\newcommand{\pykw}[1]{\textcolor{kwcol}{\textbf{#1}}}
\newcommand{\pyparam}[1]{\textcolor{paramcol}{\textbf{#1}}}
\newcommand{\pyval}[1]{\textcolor{valcol}{\textbf{#1}}}
\newcommand{\pyind}{\phantom{0000}}
\newcommand{\pycolon}{\char58\relax}

% un paramètre = un « trou » du plan ; une valeur = un trou rempli
\newcommand{\pyhole}[1]{\tikz[baseline=(h.base)]\node[draw=paramcol, dashed, rounded corners=2pt, fill=white, inner sep=1.5pt] (h) {\pyparam{#1}};}
\newcommand{\pyfilled}[1]{\tikz[baseline=(h.base)]\node[draw=valcol, rounded corners=2pt, fill=valcol!12, inner sep=1.5pt] (h) {\pyval{#1}};}

\tikzset{
  pycodeline/.style={anchor=west, inner sep=0pt, minimum height=0.4cm, font=\ttfamily\footnotesize},
  pymemval/.style={draw=black!60, fill=white, anchor=west, minimum width=1.3cm, minimum height=0.36cm, inner sep=2pt, font=\ttfamily\footnotesize\bfseries, text=valcol},
  pymemname/.style={anchor=east, font=\ttfamily\footnotesize},
}

% macros d'aide pour dessiner les slides animées (plan / appel de fonction)
% ligne n du programme : y = 6.2 - 0.42 (n-1)
\newcommand{\pyinstance}[3]{% #1 = décalage x, #2 = valeur de a, #3 = valeur de b
  \draw[draw=black!60, fill=white, rounded corners=3pt] (#1+0.1,0.95) rectangle (#1+5.0,3.7);
  \node[pycodeline] at (#1+0.3,3.4)  {minimum\_int(\pyfilled{#2}, \pyfilled{#3})};
  \node[pycodeline] at (#1+0.3,3.03) {\pyparam{a}\,=\,\pyval{#2},\ \pyparam{b}\,=\,\pyval{#3}};
  \node[pycodeline] at (#1+0.3,2.66) {\pykw{if} \pyfilled{#2} < \pyfilled{#3}\pycolon};
  \node[pycodeline] at (#1+0.3,2.29) {\pyind\pykw{return} \pyfilled{#2}};
  \node[pycodeline] at (#1+0.3,1.92) {\pykw{else}\pycolon};
  \node[pycodeline] at (#1+0.3,1.55) {\pyind\pykw{return} \pyfilled{#3}};
}
\newcommand{\pyline}[2]{\node[pycodeline] at (0,{6.2-0.42*(#1-1)}) {#2};}
\newcommand{\pyhl}[1]{\fill[yellow!45] (-0.12,{6.2-0.42*(#1-1)-0.2}) rectangle (5.4,{6.2-0.42*(#1-1)+0.2});}
\newcommand{\pyptr}[1]{\fill[red!70!black] (-0.45,{6.2-0.42*(#1-1)+0.11}) -- (-0.2,{6.2-0.42*(#1-1)}) -- (-0.45,{6.2-0.42*(#1-1)-0.11}) -- cycle;}
\newcommand{\pyfade}[1]{\fill[blueprintfill, opacity=0.8] (-0.1,{6.2-0.42*(#1-1)-0.2}) rectangle (5.38,{6.2-0.42*(#1-1)+0.2});}
\newcommand{\pymemrow}[3]{\node[pymemname] at (7.75,#1) {#2}; \node[pymemval] at (7.95,#1) {#3};}

% ------------------------------------------------------------------
% Recette « faire_frites(légumes) » : un algorithme en langage naturel
% (paramètre = trou violet, valeur = trou rempli orange, comme dans les
% slides suivantes)
% ------------------------------------------------------------------
\newcommand{\pyrecette}[1]{%
  {\ttfamily faire\_frites(#1)}\par\smallskip
  \begin{enumerate}\small\setlength{\itemsep}{0.2em}
    \item Prendre les #1
    \item Éplucher les #1
    \item Laver les #1
    \item Couper les #1 en petits dés
    \item Mettre les #1 dans un saladier
    \item Saler et assaisonner les #1
    \item Mettre les #1 au four
    \item Les sortir et les servir.
  \end{enumerate}%
}

% illustration : img/<nom>.png (ou .jpg / .jpeg) si le fichier existe, sinon un cadre de remplacement
\newcommand{\pyimage}[2]{%
  \IfFileExists{img/#1.png}{\includegraphics[width=3.2cm,height=3.2cm,keepaspectratio]{img/#1.png}}{%
  \IfFileExists{img/#1.jpg}{\includegraphics[width=3.2cm,height=3.2cm,keepaspectratio]{img/#1.jpg}}{%
  \IfFileExists{img/#1.jpeg}{\includegraphics[width=3.2cm,height=3.2cm,keepaspectratio]{img/#1.jpeg}}{%
    \tikz\node[draw=black!40, dashed, thick, rounded corners=3pt, fill=black!4, minimum width=3.2cm, minimum height=3.2cm, text width=2.8cm, align=center, text=black!55, font=\small\itshape] {image :\\#2};}}}%
}

% message d'erreur affiché à la place de l'illustration
\newcommand{\pyerreur}[1]{%
  \tikz\node[draw=red!70!black, thick, rounded corners=3pt, fill=red!8, minimum width=3.2cm, minimum height=3.2cm, text width=2.8cm, align=center, text=red!70!black, font=\small\bfseries] {#1};%
}

% ------------------------------------------------------------------
% Une fonction qui en appelle une autre : faire_frites -> couper_en_des
% ------------------------------------------------------------------
% un appel de fonction, affiché comme une ligne de code : #1 = nom, #2 = argument
\newcommand{\pycall}[2]{%
  \tikz\node[draw=black!30, fill=black!6, rounded corners=3pt, inner sep=5pt, font=\ttfamily\small] {#1(#2)};%
}

% ligne numérotée de la recette (slide « zoom ») : #1 = numéro, #2 = texte
\newcommand{\pyli}[2]{%
  \node[anchor=east, font=\footnotesize, text=black!60] at (0.4,{5.65-0.5*(#1-1)}) {#1.};
  \node[anchor=west, font=\footnotesize] at (0.5,{5.65-0.5*(#1-1)}) {#2};
}

% ligne d'un algorithme en langage naturel : #1 = niveau d'indentation, #2 = numéro de ligne, #3 = texte
\newcommand{\pynl}[3]{%
  \fill[black!45] ({0.15+0.7*#1},{5.45-0.68*(#2-1)}) circle (0.045);
  \node[anchor=west, font=\footnotesize, inner sep=2pt] at ({0.3+0.7*#1},{5.45-0.68*(#2-1)}) {#3};
}

% https://tex.stackexchange.com/questions/401884/how-do-i-change-hyperlinks-color-only
\hypersetup{
  colorlinks,
  allcolors=.,
  urlcolor=DarkBlue,
  pdftitle=Algorithmique et Programmation 1 - CM4
}


\author{}
\date{Lundi 5 octobre 2026}

\title{\includegraphics[width=0.75\textwidth]{img/logo-igm.png} \\[1em] Algorithmique et Programmation 1}
\subtitle{Approfondissements sur les fonctions}
\institute{L1 Mathématiques - L1 Informatique \\ Semestre 1}

\begin{document}

\maketitle

% Objectifs : Fonctions complexes: mémoire, prototypes/signature Appels : évaluation, retour, pile. Terminal. Exemple : recherche du minimum dans une liste

\section{Approfondissements : les fonctions}

\begin{frame}{Fonction}
    En informatique, une fonction est :
    \begin{itemize}
    \item Un morceau de programme
    \item Portant en général \emph{un nom}
    \item Prenant un ou plusieurs \emph{paramètres} (ou zéro)
    \item Renvoyant un résultat (la plupart du temps)
    \end{itemize}
\end{frame}

\begin{frame}{Anatomie d'une fonction}
    \begin{tikzpicture}
      % Rectangle pour représenter la fonction
      \node [draw, rectangle, text width=2cm, minimum height=1.5cm, align=center] (function) at (0,0) {Fonction};

      % Chemin pour représenter l'entrée (paramètres)
      \draw[-latex] (-5,0) -- (-1.2,0) node[midway, above, label={below: (paramètres)}] {Entrée};

      % Chemin pour représenter la sortie (valeur de retour)
      \draw[-latex] (1.2,0) -- (5,0) node[midway, above, label= {below: (valeur de retour)}] {Sortie};

      % Flèche d'affichage (zig-zag)
      \onslide<2->{\draw [decorate, -latex, decoration={snake, segment length=5mm, amplitude=1mm, post length=1mm}] (0,0.5) -- (0,2.5) node[midway, left] {Affichages};}

      % Flèche d'effets de bord (zig-zag)
      % \onslide<3->{\draw [decorate, -latex, decoration={snake, segment length=5mm, amplitude=1mm, post length=1mm}] (0,-0.5) -- (0,-2.5) node[midway, left, label={right: (aka: Effets secondaires)}] {Effets de bord};}
    \end{tikzpicture}
\end{frame}

% ==================================================================
% Recette : faire_frites(légumes), avec un paramètre qu'on remplace
% ==================================================================

\begin{frame}{Un algorithme de cuisine}
  \begin{columns}[c]
    \begin{column}{0.68\textwidth}
      \only<1>{\pyrecette{\pyhole{légumes}}}%
      \only<2>{\pyrecette{\pyfilled{patates}}}%
      \only<3>{\pyrecette{\pyfilled{patates douces}}}%
      \only<4>{\pyrecette{\pyfilled{courgettes}}}%
      \only<5>{\pyrecette{\pyfilled{yaourts}}}%
    \end{column}
    \begin{column}{0.30\textwidth}
      \centering
      \only<2>{\pyimage{frites_patates}{frites de patates}}%
      \only<3>{\pyimage{frites_patates_douces}{frites de patates douces}}%
      \only<4>{\pyimage{frites_courgettes}{frites de courgettes}}%
      \only<5>{\pyerreur{Erreur : Impossible d'éplucher un yaourt}}%
    \end{column}
  \end{columns}
  \vspace{0.2cm}
  \begin{center}
    \parbox[c][0.9cm][c]{\textwidth}{\centering
      \only<2>{\pycall{faire\_frites}{\pyfilled{patates}}}%
      \only<3>{\pycall{faire\_frites}{\pyfilled{patates douces}}}%
      \only<4>{\pycall{faire\_frites}{\pyfilled{courgettes}}}%
      \only<5>{\pycall{faire\_frites}{\pyfilled{yaourts}}}%
    }
  \end{center}
\end{frame}

\begin{frame}[fragile]{Fonction à paramètres et résultats}
  Calculer le \emph{minimum} de deux entiers (dans Thonny) :
  \pause
\begin{minted}{python}
def minimum_int(a, b):
    if a < b:
        return a
    else:
        return b
\end{minted}
\end{frame}

% ==================================================================
% Slide A : une fonction, c'est un plan (paramètres = trous à remplir)
% ==================================================================

\begin{frame}{Une fonction, c'est un plan}
  \centering
  \begin{tikzpicture}[x=1cm,y=1cm]
    \useasboundingbox (0,0.2) rectangle (10.6,6.7);

    % --- le plan : la définition, avec des « trous » ---
    \draw[blueprintline, dashed, thick, fill=blueprintfill, rounded corners=3pt] (2.85,4.25) rectangle (7.75,6.3);
    \node[anchor=south west, font=\scriptsize\bfseries, text=blueprintline] at (2.85,6.32) {LE PLAN (la définition)};
    \node[pycodeline] at (3.05,6.02) {\pykw{def} minimum\_int(\pyhole{a}, \pyhole{b})\pycolon};
    \node[pycodeline] at (3.05,5.62) {\pyind\pykw{if} \pyhole{a} < \pyhole{b}\pycolon};
    \node[pycodeline] at (3.05,5.22) {\pyind\pyind\pykw{return} \pyhole{a}};
    \node[pycodeline] at (3.05,4.82) {\pyind\pykw{else}\pycolon};
    \node[pycodeline] at (3.05,4.42) {\pyind\pyind\pykw{return} \pyhole{b}};

    \only<1>{
      \node[anchor=west, text=paramcol, font=\small\itshape, align=left] (lblpar) at (8.1,6.02) {paramètres :\\les « trous »};
      \draw[->, paramcol, thick] (lblpar.west) -- (7.05,6.02);
    }

    % --- premier appel ---
    \only<2->{
      \draw[->, thick, blueprintline] (3.6,4.23) -- (2.6,3.72);
      \node[anchor=east, font=\scriptsize\itshape, text=blueprintline] at (2.2,4.0) {on remplit les trous};
      \pyinstance{0}{58}{78}
      \node[anchor=south east, font=\scriptsize\bfseries, text=black!55] at (4.95,3.72) {APPEL 1};
    }
    \only<3->{
      \fill[white, opacity=0.7] (0.2,1.35) rectangle (4.95,2.1);
      \node[anchor=west, font=\scriptsize\itshape, text=black!60] at (3.7,2.66) {vrai};
      \node[pycodeline, text=valcol] at (0.3,1.15) {$\Rightarrow$ \textbf{renvoie} \pyval{58}};
    }

    % --- second appel, même plan ---
    \only<4->{
      \draw[->, thick, blueprintline] (7.0,4.23) -- (8.05,3.72);
      \pyinstance{5.5}{12}{3}
      \node[anchor=south east, font=\scriptsize\bfseries, text=black!55] at (10.45,3.72) {APPEL 2};
    }
    \only<5->{
      \fill[white, opacity=0.7] (5.7,2.1) rectangle (10.45,2.48);
      \node[anchor=west, font=\scriptsize\itshape, text=black!60] at (9.2,2.66) {faux};
      \node[pycodeline, text=valcol] at (5.8,1.15) {$\Rightarrow$ \textbf{renvoie} \pyval{3}};
    }

    % --- légendes ---
    \node[font=\small, align=center, text width=10.4cm] at (5.3,0.5) {%
      \only<1>{Le plan est écrit \emph{une seule fois}. Ses paramètres sont des « trous » à remplir.}%
      \only<2>{À chaque appel, les \emph{arguments} remplissent les trous.}%
      \only<3>{On exécute le plan avec ces valeurs : \texttt{58 < 78} est vrai.}%
      \only<4>{Même plan, autres valeurs.}%
      \only<5>{\texttt{12 < 3} est faux : on renvoie 3. Un seul plan, des résultats différents.}%
    };
  \end{tikzpicture}
\end{frame}

% ==================================================================
% Slide B : que se passe-t-il lors d'un appel ? (style Python Tutor)
% ==================================================================
% ligne n du programme : y = 6.2 - 0.42 (n-1)

\begin{frame}{Que se passe-t-il lors d'un appel ?}
  \centering
  \begin{tikzpicture}[x=1cm,y=1cm]
    \useasboundingbox (-0.9,-0.1) rectangle (9.85,6.65);

    % --- le plan (lignes 1 à 5) en fond ---
    \fill[blueprintfill] (-0.12,4.32) rectangle (5.4,6.4);
    \draw[blueprintline, dashed] (-0.12,4.32) rectangle (5.4,6.4);

    % --- ligne courante et pointeur ---
    \only<1>{\pyhl{1}\pyptr{1}}
    \only<2>{\pyhl{7}\pyhl{8}\pyptr{8}}
    \only<3-4>{\pyhl{9}\pyptr{9}}
    \only<5>{\pyhl{1}\pyptr{1}}
    \only<6>{\pyhl{2}\pyptr{2}}
    \only<7>{\pyhl{3}\pyptr{3}}
    \only<8>{\pyhl{9}\pyptr{9}}
    \only<9>{\pyhl{10}\pyptr{10}}

    % --- le programme ---
    \pyline{1}{\pykw{def} minimum\_int(\pyhole{a}, \pyhole{b})\pycolon}
    \only<1-5,8->{\pyline{2}{\pyind\pykw{if} \pyhole{a} < \pyhole{b}\pycolon}}
    \only<6-7>{\pyline{2}{\pyind\pykw{if} \pyfilled{58} < \pyfilled{78}\pycolon}}
    \only<1-5,8->{\pyline{3}{\pyind\pyind\pykw{return} \pyhole{a}}}
    \only<6-7>{\pyline{3}{\pyind\pyind\pykw{return} \pyfilled{58}}}
    \pyline{4}{\pyind\pykw{else}\pycolon}
    \only<1-5,8->{\pyline{5}{\pyind\pyind\pykw{return} \pyhole{b}}}
    \only<6-7>{\pyline{5}{\pyind\pyind\pykw{return} \pyfilled{78}}}
    \only<7>{\pyfade{4}\pyfade{5}}

    \pyline{7}{x = \pyval{58}}
    \pyline{8}{y = \pyval{78}}
    \only<1-2,9>{\pyline{9}{mon\_min = minimum\_int(x, y)}}
    \only<3>{\pyline{9}{mon\_min = minimum\_int(\pyval{x}, \pyval{y})}}
    \only<4-7>{\pyline{9}{mon\_min = minimum\_int(\pyfilled{58}, \pyfilled{78})}}
    \only<8>{\pyline{9}{mon\_min = \pyfilled{58}}}
    \pyline{10}{print(mon\_min)}
    \node[anchor=north east, font=\scriptsize\itshape, text=blueprintline] at (5.37,6.38) {plan};

    % --- appel et retour (flèches pointillées dans la marge) ---
    \only<5>{\draw[->, dashed, thick, valcol] (-0.55,{6.2-0.42*8}) .. controls (-1.05,{6.2-0.42*8}) and (-1.05,6.2) .. (-0.55,6.2);}
    \only<8>{\draw[->, dashed, thick, valcol] (-0.55,{6.2-0.42*2}) .. controls (-1.05,{6.2-0.42*2}) and (-1.05,{6.2-0.42*8}) .. (-0.55,{6.2-0.42*8});}

    % --- mémoire : espace global ---
    \draw[rounded corners=3pt, draw=black!50, fill=black!4] (5.75,4.05) rectangle (9.85,6.6);
    \node[anchor=north west, font=\scriptsize\bfseries, text=black!70] at (5.85,6.56) {Espace global};
    \node[pymemname] at (7.75,5.75) {minimum\_int};
    \node[pymemval, draw=blueprintline, dashed, fill=blueprintfill, text=blueprintline, font=\scriptsize\itshape] at (7.95,5.75) {plan};
    \only<2->{\pymemrow{5.3}{x}{58}\pymemrow{4.85}{y}{78}}
    \only<8->{\pymemrow{4.4}{mon\_min}{58}}
    \only<3>{\draw[valcol, thick, rounded corners=2pt] (5.82,5.09) rectangle (9.4,5.51);
             \draw[valcol, thick, rounded corners=2pt] (5.82,4.64) rectangle (9.4,5.06);}

    % --- mémoire : espace local de l'appel ---
    \only<5-7>{
      \draw[rounded corners=3pt, draw=paramcol!70, fill=paramcol!6] (5.75,1.85) rectangle (9.85,3.85);
      \node[anchor=north west, font=\scriptsize\bfseries, text=paramcol] at (5.85,3.81) {Espace local (appel)};
      \pymemrow{3.2}{\pyparam{a}}{58}
      \pymemrow{2.75}{\pyparam{b}}{78}
    }
    \only<7>{\pymemrow{2.3}{\textbf{renvoie}}{58}}

    % --- affichage ---
    \only<9>{
      \draw[rounded corners=3pt, draw=black!50, fill=black!8] (5.75,0.85) rectangle (9.85,1.85);
      \node[anchor=north west, font=\scriptsize\bfseries, text=black!70] at (5.85,1.81) {Affichage};
      \node[anchor=west, font=\ttfamily\footnotesize] at (5.95,1.15) {58};
    }

    % --- légendes ---
    \node[font=\small, align=center, text width=10.2cm] at (4.5,0.4) {%
      \only<1>{La définition range le plan en mémoire. Le corps n'est pas exécuté.}%
      \only<2>{Le programme principal commence : \texttt{x} et \texttt{y} sont créées.}%
      \only<3>{Appel : on évalue d'abord les arguments, ici les variables \texttt{x} et \texttt{y}.}%
      \only<4>{Les variables \texttt{x} et \texttt{y} sont remplacées par leurs valeurs : \pyval{58} et \pyval{78}.}%
      \only<5>{Un espace local est créé : le paramètre \pyparam{a} reçoit \pyval{58}, \pyparam{b} reçoit \pyval{78}.}%
      \only<6>{Dans le corps, \pyparam{a} et \pyparam{b} sont remplacés par leurs valeurs : \texttt{58 < 78} est vrai.}%
      \only<7>{\texttt{return} : la fonction s'arrête et renvoie \pyval{58}.}%
      \only<8>{L'appel est remplacé par la valeur renvoyée, affectée à \texttt{mon\_min}. L'espace local disparaît.}%
      \only<9>{\texttt{print} affiche la valeur de \texttt{mon\_min}.}%
    };
  \end{tikzpicture}
\end{frame}

% ==================================================================
% Une fonction peut en appeler une autre : zoom sur couper_en_des
% ==================================================================

\begin{frame}{Une fonction peut en appeler une autre}
  \centering
  \begin{tikzpicture}[x=1cm,y=1cm]
    \useasboundingbox (0,0.6) rectangle (10.6,6.6);

    % le zoom : cette étape est elle-même un appel de fonction (dessiné derrière le texte)
    \only<2->{
      \fill[yellow!30] (5.55,4.37) -- (6.15,5.2) -- (6.15,3.1) -- (5.55,3.93) -- cycle;
    }

    % l'étape sur laquelle on zoome
    \fill[yellow!45] (-0.05,3.93) rectangle (5.55,4.37);

    % la recette pour les patates (le contexte, en petit)
    \node[anchor=west, font=\ttfamily\footnotesize] at (0,6.2) {faire\_frites(\pyfilled{patates})};
    \pyli{1}{Prendre les \pyfilled{patates}}
    \pyli{2}{Éplucher les \pyfilled{patates}}
    \pyli{3}{Laver les \pyfilled{patates}}
    \pyli{4}{Couper les \pyfilled{patates} en petits dés}
    \pyli{5}{Mettre les \pyfilled{patates} dans un saladier}
    \pyli{6}{Saler et assaisonner les \pyfilled{patates}}
    \pyli{7}{Mettre les \pyfilled{patates} au four}
    \pyli{8}{Les sortir et les servir.}

    \only<2->{
      \draw[rounded corners=4pt, thick, draw=black!60, fill=white] (6.15,3.1) rectangle (11.0,5.2);
      \node[font=\scriptsize\itshape, text=black!55] at (8.575,4.85) {appel de fonction};
      \node at (8.575,4.1) {\pycall{couper\_en\_des}{\pyfilled{patates}}};
    }
  \end{tikzpicture}
\end{frame}

\begin{frame}{Zoom sur \texttt{couper\_en\_des}}
  \centering
  \begin{tikzpicture}[x=1cm,y=1cm]
    \useasboundingbox (0,0.6) rectangle (10.6,6.6);

    \node[anchor=west, font=\ttfamily] at (0,6.2) {couper\_en\_des(\pyhole{legumes})};

    \only<1->{\pynl{0}{1}{\pykw{Pour chaque} légume \pykw{dans} \pyhole{\texttt{legumes}} :}}
    \only<2->{\pynl{1}{2}{\pykw{Tant que} \texttt{largeur(legume)} > 2~cm \pykw{ou} \texttt{longueur(legume)} > 2~cm :}}
    \only<3->{\pynl{2}{3}{\pykw{Si} \texttt{largeur(legume)} > \texttt{longueur(legume)} :}}
    \only<4->{\pynl{3}{4}{Couper le légume dans la largeur}}
    \only<5->{\pynl{2}{5}{\pykw{Sinon} :}}
    \only<6->{\pynl{3}{6}{Couper le légume dans la longueur}}
    \only<7->{\pynl{0}{7}{\pykw{Renvoyer} les dés de légume}}
  \end{tikzpicture}
\end{frame}

\section{Exemples}

\begin{frame}[fragile]{Fonction à paramètres et résultats}
  Calculer le \emph{minimum} d'une liste non-vide (dans Thonny) :
  \pause
  \onslide<2>
\begin{minted}{python}
def minimum_liste(lst):
    m = lst[0] # /!\ Erreur si lst est vide
    for elt in lst:
        if elt < m:
        m = elt
    return m
\end{minted}
\end{frame}

\begin{frame}[fragile]{Fonction sans valeur de retour}
  \begin{columns}
    \begin{column}{0.7\textwidth}
  Dessiner un carré fait du caractère \mintinline{python}{carac} (dans Thonny) :
\end{column}
\begin{column}{0.1\textwidth}
\begin{verbatim}
*****
*****
*****
*****
*****
\end{verbatim}
\end{column}
\end{columns}
\pause
\vspace{-1em}
\begin{overprint}
  \onslide<2>
\begin{minted}{python}
def dessine_carre(n, carac):
    i = 0
    while i < n:
        j = 0
        ligne =""
        while j < n:
            ligne += carac
            j += 1
        print(ligne)
        i += 1
    return  # Facultatif : ne renvoie rien
\end{minted}
  \onslide<3>
  Mieux :
\begin{minted}{python}
def dessine_carre(n, carac):
    for i in range(0, n):
        ligne = ""
        for j in range(0, n):
            ligne += carac
        print(ligne)
\end{minted}
\end{overprint}
\end{frame}

\begin{frame}
  \frametitle{Composition de fonctions}

  \begin{block}{En maths}
    $$\forall x \in X, g \circ f(x) = g(f(x))$$
  \end{block}

  \begin{block}{En informatique}
    \vspace{0pt}
    On peut appeler une fonction dans une fonction ! \newline (et ainsi de suite)
  \end{block}
\end{frame}

\begin{frame}[fragile]{Composition de fonctions : minimum d'une liste}
  Calculer le \emph{minimum} d'une liste non-vide (dans Thonny) :
  \pause
\begin{overprint}
  \onslide<1>
\begin{minted}{python}
def minimum_liste(lst):
    m = lst[0] # /!\ Erreur si lst est vide
    for elt in lst:
        if elt < m:
        m = elt
    return m
\end{minted}
  \onslide<2>
\begin{minted}{python}
def minimum_liste(lst):
    m = lst[0] # /!\ Erreur si lst est vide
    for elt in lst:
        m = minimum_int(m, elt)
    return m
\end{minted}
  \end{overprint}
\end{frame}

\begin{frame}[fragile]{Erreur fréquente : confusion \begin{tabular}{l}paramètre / saisie \\ retour / affichage\end{tabular}}

    \vspace{-0.25em}
  \begin{columns}
    \begin{column}{0.5\textwidth}
  \begin{block}{Paramètre / saisie}
    \vspace{-0.25em}
\begin{minted}{python}
def minimum_int(a,b):
    a = int(input()) # NON !
    b = int(input()) # NON !
    if a <= b:
        return a
    else:
        return b
\end{minted}
  \end{block}
\pause
    \vspace{-0.25em}
  \begin{block}{Retour / affichage}
    \vspace{-0.25em}
\begin{minted}{python}
def minimum_int(a,b):
    if a <= b:
        print(a) # NON !
    else:
        print(b) # NON !
\end{minted}
  \end{block}
    \end{column}
    \pause
    \begin{column}{0.5\textwidth}
      Ce que ça donne dans Thonny :

      \texttt{minimum\_liste\_input\_print.py}
    \end{column}
    \end{columns}
\end{frame}

\begin{frame}
  \frametitle{Fonctions et espaces de nom}
  \begin{itemize}
  \item Paramètres et variables d'une fonction : \alert{indépendantes} des autres variables du programme
  \item N'existent plus une fois l'exécution de la fonction terminée
  \item On les appelle des variables \emph{locales}
\end{itemize}

$\Rightarrow{}$ On peut renommer les variables d'une fonction.

$\rightarrow{}$ Démonstration sur Thonny (\texttt{minimum.py}).

\end{frame}

\begin{frame}
  \frametitle{(Contre-)exemple : intervertir des variables}
  Dans Thonny : \texttt{echange.py}
  \pause
  \begin{itemize}
  \item changer les valeurs de \mintinline{python}{a} et \mintinline{python}{b} dans la fonction \emph{n'a pas d'effet} sur \mintinline{python}{x} et \mintinline{python}{y} dans le programme principal !
  \item la variable \mintinline{python}{temp} n'existe plus après l'exécution de la fonction
  \end{itemize}
\end{frame}

\begin{frame}
  \frametitle{Sous le capot : sémantique d'un appel}

  \begin{block}{Espace de noms}
    \vspace{0pt}
    Ensemble de noms (variables, fonctions) défini à un certain point d'un programme
  \end{block}

  \MVRightarrow{} lors d'un appel de fonction, création d'un espace de nom \emph{local} :
  \begin{itemize}
  \item Paramètres $\rightarrow$ leur valeur lors de l'appel
  \item Variables locales
  \end{itemize}

  \begin{block}{Exemple}
    \vspace{0pt}
    \texttt{minimum\_liste.py} dans \href{https://pythontutor.com/}{Python Tutor} et Thonny.
  \end{block}

\end{frame}

\begin{frame}
  \frametitle{Espaces de noms lors de l'exécution}
  Les espaces de noms sont imbriqués :
  \begin{itemize}
  \item espace de noms prédéfini (\mintinline{python}{print}, \mintinline{python}{input}, \dots)
  \item espace de noms global : définis dans le programme principal (\emph{variables globales})
  \item espaces de noms locaux imbriqués dans l'\emph{ordre chronologique} des appels
    \begin{itemize}
    \item paramètres de l'appel
    \item variables locales à la fonction
    \end{itemize}
  \end{itemize}
\end{frame}

\begin{frame}[fragile]
  \frametitle{Pile d'appels}
  L'empilement des espaces de noms obéit à une politique de \emph{pile}

  \begin{itemize}
  \item sommet : appel en cours
  \item en-dessous : appels précédents
  \item (presque) tout en bas : espace de nom global
  \end{itemize}

  \begin{columns}
    \begin{column}{0.5\textwidth}
\begin{minted}{python}
def f(x):
    return g(x) + 1
def g(x):
    return h(x) + 2
def h(x):
    return x + 3
\end{minted}
    \end{column}
    \begin{column}{0.4\textwidth}
      \alt<1-4>{%
      \begin{tikzpicture}
        % Stack rectangle
        \draw (0,0) rectangle (2,3);
        % Stack top
        \node at (1,3.3) {Sommet};
        % Stack elements
        \pause
        \draw (0.2,0.5) rectangle (1.8,0.8) node[midway] {f};
        \pause
        \draw (0.2,0.9) rectangle (1.8,1.2) node[midway] {g};
        \pause
        \draw (0.2,1.3) rectangle (1.8,1.6) node[midway] {h};
        % Stack direction
        \draw[->,ultra thick] (2.5,0) -- (2.5,3);
      \end{tikzpicture}%
    }{%
      \begin{tikzpicture}[yscale=-1]
        % Stack rectangle
        \draw (0,0) rectangle (2,3);
        % Stack top
        \node at (1,3.3) {Sommet};
        % Stack elements
        \pause
        \draw (0.2,0.5) rectangle (1.8,0.8) node[midway] {f};
        \pause
        \draw (0.2,0.9) rectangle (1.8,1.2) node[midway] {g};
        \pause
        \draw (0.2,1.3) rectangle (1.8,1.6) node[midway] {h};
        % Stack direction
        \draw[->,ultra thick] (2.5,0) -- (2.5,3);
      \end{tikzpicture}%
    }
    \end{column}
  \end{columns}

  \medskip

  \onslide<5>{\textbf{Attention !} Dans Python Tutor, le plus récent est en bas.}
\end{frame}

\begin{frame}
  \frametitle{Pile d'appels}

  \begin{block}{Quand un nouvel appel commence}
    \begin{itemize}
    \item l'exécution de la fonction en cours s'interrompt
    \item un nouvel espace de noms local est créé
    \item l'exécution de la fonction appelée commence
    \end{itemize}
  \end{block}


  \begin{block}{Quand l'appel en cours se termine}
    \begin{itemize}
    \item son espace de nom est supprimé de la pile
    \item l'exécution de l'appel précédent reprend
    \end{itemize}

  \end{block}

\end{frame}
\begin{frame}{Portée des variables}

  \begin{block}{Accès à la valeur d'une variable}
  \begin{itemize}
  \item possible pour n'importe quel nom défini dans un des espaces de noms antérieurs
  \item si plusieurs espaces contiennent le même nom, c'est le plus récent qui est sélectionné
  \item la plupart du temps, on n'a besoin que des variables locales

    (quitte à ajouter des paramètres)
  \end{itemize}
\end{block}

\pause

\begin{block}{Affectation}
  \vspace{0pt}
  Seulement les variables \alert{locales} de la fonction courante !

  (ou globales \emph{si on est dans le programme principal})
\end{block}

\end{frame}

\begin{frame}[fragile]{Signature d'une fonction}

Considérons la fonction :
\begin{minted}{python}
def f(p_1, ..., p_n):
...
return expr
\end{minted}

Les noms \mintinline{python}{p_1} à \mintinline{python}{p_n} sont appelés paramètres (formels).

En général, ils d'un certain type. On peut le préciser en explicitant la \emph{signature} de la fonction :
\begin{minted}{python}
def f(p_1: type_1, ..., p_n: type_n) -> type_retour:
...
return expr
\end{minted}
\begin{exampleblock}{Exemple}
\begin{minted}{python}
def minimum_int(a: int, b: int) -> int:
\end{minted}
\end{exampleblock}
\end{frame}

\begin{frame}[fragile]{Déroulement détaillé d'un appel}
      Considérons maintenant l'appel \mintinline{python}{f(e_1, ..., e_n)} :
      \begin{itemize}
      \item les valeurs \mintinline{python}{v_1} à \mintinline{python}{v_n} sont appelées paramètres effectifs \newline (ou arguments)
      \item création d'un espace de noms local contenant \mintinline{python}{p_1} à \mintinline{python}{p_n} au sommet de la pile d'appels
      \item chaque expression \mintinline{python}{e_i} est évaluée en une valeur \mintinline{python}{v_i} et affectée à la variable \mintinline{python}{p_i}
      \item exécution du corps de la fonction dans l'espace de noms local
        \pause
      \item si la fonction exécute l'instruction \mintinline{python}{return expr} ou atteint la fin de son bloc d'instructions :
        \begin{itemize}
        \item l'espace de noms local est détruit
        \item l'expression appelante \mintinline{python}{f(e_1, ..., e_n)} prend la valeur de \mintinline{python}{expr} (respectivement \mintinline{python}{None})
        \item reprise du programme principal dans l'espace global
        \end{itemize}
      \end{itemize}
\end{frame}

\section{Bonnes pratiques de programmation}

\begin{frame}[fragile]{Documentation et test de fonctions}
  \begin{block}{Chaînes de documentation (\mintinline{python}{docstring})}
    \vspace{0pt}

    Bonne pratique : indiquer par un commentaire
    \vspace{-0.5em}
    \begin{itemize}
    \item à quoi sert une fonction
     \item ce que représentent ses paramètres et leur type
     \item ce que représente sa valeur de retour
     \item d'éventuels effets ou causes secondaires
    \end{itemize}
  \end{block}

\begin{minted}{python}
def triple(n):
  """
  Calcule le triple du nombre n (int ou float)
  ou la répétition trois fois de la chaîne n.
  """
  return n * 3
\end{minted}
\end{frame}
\begin{frame}[fragile]{Documentation et test de fonctions}
  \begin{itemize}
    \item On peut accéder à la chaîne de documentation d'une fonction en tapant \mintinline{python}{help(nom de la fonction)} dans l'interpréteur
  \item Cela fonctionne aussi pour les fonctions prédéfinies ou issues de modules
  \end{itemize}
\end{frame}
\begin{frame}[fragile]{Tests intégrés à la documentation (\mintinline{python}{doctest})}
  \begin{itemize}
  \item Toute fonction \emph{doit être testée immédiatement} pour s'assurer qu'elle fonctionne.
  \item On peut intégrer les tests à sa documentation.
  \end{itemize}
\begin{minted}{python}
def triple(n):
  """
  Calcule le triple du nombre n (int ou float)
  ou la répétition trois fois de la chaîne n.

  >>> triple(3)
  9
  >>> triple(9.0)
  27.0
  """
  return n * 3
\end{minted}
\end{frame}
\begin{frame}[fragile,fragile]{Tests intégrés à la documentation (\mintinline{python}{doctest})}
  Il existe des outils qui permettent de lancer automatiquement tous les tests présents dans la documentation, et de vérifier qu'ils produisent les résultats annoncés.

  Par exemple, à la fin d'un programme, on peut écrire le code suivant pour lancer systématiquement tous les tests présents dans le fichier :
\begin{minted}{python}
  import doctest
  doctest.testmod()
\end{minted}
\end{frame}

\section{Extra}

\begin{frame}{Responsabilité}

  En général, une fonction est \emph{responsable} d'\alert{une seule chose}.

  Sinon, on fait \emph{plusieurs} fonctions.

      Ce que ça donne dans Thonny :

      \texttt{minimum\_liste\_responsabilite.py}
\end{frame}

\begin{frame}[fragile]{Un autre exemple de composition de fonctions : dessiner un carré}
  \begin{columns}
    \begin{column}{0.1\textwidth}
\begin{verbatim}
*****
*****
*****
*****
*****
\end{verbatim}
    \end{column}
    \begin{column}{0.7\textwidth}
      Carré de côté \mintinline{python}{n} = \mintinline{python}{n} lignes de longueur \mintinline{python}{n}
    \end{column}
\end{columns}
\bigskip
\pause
On peut donc décomposer ainsi :
  \begin{itemize}
  \item
\begin{minted}{python}
def dessine_ligne(n,carac):
    """Dessine *une ligne* de longueur n
       composée du caractère carac."""
\end{minted}
  \item
\begin{minted}{python}
def dessine_carre(n,carac):
    """Dessine *un carré* de longueur n
       composé du caractère carac."""
\end{minted}
  \end{itemize}
  % \MVRightarrow{} Ce que ça donne dans Thonny
\end{frame}
\begin{frame}[fragile]{Dessiner un carré : variante}
\begin{minted}{python}
def dessine_ligne(n, carac):
 """dessine une ligne de longueur n
    composée du caractère carac"""
    for j in range(0, n):
        print(caractere, end = '')
    print()
    return
\end{minted}
\end{frame}
\begin{frame}[fragile]{Dessiner un carré : variante}
\begin{minted}{python}
def dessine_carre(n, caractere):
 """dessine un carré de longueur n
    composé du caractère carac"""
    for i in range(0, n):
        dessine_ligne(n,caractere)
    return
\end{minted}
\end{frame}


\end{document}

%%% Local Variables:
%%% TeX-engine: xetex
%%% TeX-command-extra-options: "-shell-escape"
%%% End:
